\chapter{Validating the measurement}
\label{app:validation}

This appendix holds the parts of the evaluation that establish what the measuring
instruments can and cannot see. Nothing here is a result of the analysis; it is the
record that the figures of Chapter~\ref{ch:results} were produced by instruments whose
own failures were looked for and found.

\section{Nodes the two static readers disagreed about}
\label{app:readers}

Table~\ref{tab:coverage} also reports the other direction of the same comparison, which
the differential test produces and which an earlier version of the results chapter left out. Every
figure in this section is the first round's, measured with the parser as it stood then;
Section~\ref{sec:res-round2} repeats the comparison with both tools corrected and the
numbers move a long way. As the first round measured it, the
parser holds no declared screen for 197 of the 338 nodes in the corpus, so node extraction
is as inaccurate as it is incomplete. For k8s dashboard the figure is almost entirely the
parser's own limit, 37 of its 38 unmatched nodes being the lazily loaded entries it cannot
follow. Mifos is not explained that way: it has one blind spot and 120 unmatched nodes.
Table~\ref{tab:node-audit} decomposes all 197 by the kinds of route the parser found at
the node's path, joined on the parser's own normalisation.

\begin{table}[t]\footnotesize
\caption{Why each unmatched node has no declared screen. A node is placed in the first
row that applies, and the categories are exclusive.}
\label{tab:node-audit}
\TabNodeAudit
\end{table}

Three groups come out of it. For 43 nodes the parser could not follow the route that would
have declared them, so those are the oracle's limit and not the extractor's, and 37 of the
43 are k8s dashboard. For 11 nodes the parser did find a route, but one it classifies as a
path prefix or a redirect. Section~\ref{sec:model} admits only component-bearing routes
into the screen set, so the implemented node set contains 11 elements the model excludes:
one in realworld, two in alfresco, six in bitwarden and two in mifos. The clearest case is
\texttt{/profile} in realworld, whose route declares \texttt{loadChildren} and no
component. These nodes make $\mathit{screen}(r)$ answerable for a route the model treats
as a path prefix only.

The remaining \RoundTwoFirstRoundExtra{} are the substantial group and have no declared route
of any kind at the node's path. They are concentrated in mifos, which has 117 of them against a single parser
blind spot, and in bitwarden, which has 23. Either the extractor emits nodes for paths it
assembles differently from the parser, or the parser misses route tables it does not report
as blind spots. The two methods disagree here without either being able to adjudicate, and
at the close of the first round it was the largest single piece of unexplained behaviour
the evaluation had found, sitting underneath every figure that counts nodes.

The second round resolved it, and the answer was the FIRST of the two explanations almost
throughout --- the opposite of what an earlier draft concluded.
\RoundTwoFirstRoundTruncated{} of the \RoundTwoFirstRoundExtra{} sit at an address that is the
tail of a route the parser had already declared in full: the extractor had reached the table
without its parent prefix and minted \path{/premium} where the parser declared
\path{/settings/subscription/premium}. That is the extractor assembling the path differently,
not the parser missing a table. The other \RoundTwoFirstRoundSynthetic{} are the catch-all
each of those five applications declares, \texttt{path: '**'} with a component, which the
extractor mints at the path \texttt{*} while the oracle's policy keeps it out of the screen
set: there the two tools disagree by convention, not over a table either of them missed. Which way the group fell matters beyond the bookkeeping, because it is the
difference between a measuring instrument that was wrong and a tool that was: the repairs of
Section~\ref{sec:res-round2} are accordingly almost all in the extractor. The figures above are the first round's and are left standing as the record of
what that round saw; Section~\ref{sec:res-round2} repeats the comparison with both tools
corrected, and the corpus is then left with \RoundTwoExtraNodes{} nodes that no declared
screen matches, of which \RoundTwoExtraNodesWildcard{} sit at the path \texttt{*}, where
the extractor mints the catch-all the application declares and the oracle counts no screen. The parser had its own share of the
problem, counted in Section~\ref{sec:res-round2}: a route array behind a module it refused to
follow, a path assembled from a constant it did not resolve, a lazily loaded table walked
without its parent. It was the smaller share. Reporting the disagreement
rather than attributing it was the right call at the time, and working it is what produced
most of Section~\ref{sec:res-round2}.

So the baseline recovers a small fraction of the navigation, and loses it at two levels: some screens never become nodes, and most navigation sites of the screens that do are not turned into edges.



On the instrument the figures of Chapter~\ref{ch:results} are measured on, this
disagreement is gone: the parser declares a screen for every node the extractor mints
except the wildcard catch-all route, which it does not count as a screen, and the runtime
method confirms that the framework mounts a component at every address the extractor
claims. What this section records is therefore not a standing defect but the state the
differential test was in when it was the only check available, and the reason a third,
independent reader was added.

\section{How the declared-screen set moved}
\label{app:setmoved}


Table~\ref{tab:round2nodes} gives node coverage before and after. \RoundTwoFullNodes{} of
the six applications now hold a node for every declared screen. The denominator of those ratios
deserves more than a column of its own, because on this corpus the declared-screen \emph{count}
turns out to be a poor summary of the declared-screen \emph{set}: between the oracle as it stood
before this round and the oracle reported here, \RoundTwoSetMoved{} of the six applications'
declared sets changed, and \RoundTwoSetMovedCountSame{} of those changed while the count stood
still. A moving denominator is how a coverage figure is rigged without anyone intending it, and a
denominator that holds its size while every address inside it changes is how the same thing
happens without even the column moving. Each of the five is accounted for below, from the
generated subtraction rather than from a sentence.

ghostfolio is the clearest case and the one that would have passed unnoticed. Its count did not
move at all: \RoundTwoEarlierOracleGhostfolioDeclared{} declared screens before,
\RoundTwoEarlierOracleGhostfolioNow{} after. Every single address in it changed.
\RoundTwoEarlierOracleGhostfolioInside{} paths left the set and
\RoundTwoEarlierOracleGhostfolioAdded{} arrived, because the earlier oracle recorded the verbatim
expression a route's path was written as rather than the address it resolves to ---
\path{/internalRoutes.account.path} is what it stored, where the application mounts
\path{/about} and its siblings. A reader comparing the two counts would have concluded that
nothing about ghostfolio's route set had changed, and would have been wrong about all of it.

k8s dashboard carries the largest move, and it was missing from this account until the
provenance of the comparison was fixed. It went from \RoundTwoEarlierOracleKdashDeclared{}
declared screens to \RoundTwoEarlierOracleKdashNow{}, with
\RoundTwoEarlierOracleKdashAdded{} arriving and none leaving. The earlier oracle reported 37
blind spots for this application, every one of them a \texttt{loadChildren} target it could not
resolve, so it declared the two screens its two eagerly registered tables carry ---
\path{/login} and the application root --- and nothing beneath any lazy module; the oracle
reported here declares the whole tree with no blind spot. That the earlier figure was a near-vacuous denominator, and that
coverage against it was therefore 100\,\% of almost nothing, is a fact about the measuring
instrument and not about the application.

bitwarden went from \RoundTwoEarlierOracleBitwardenDeclared{} declared screens to
\RoundTwoEarlierOracleBitwardenNow{}. \RoundTwoEarlierOracleBitwardenStorybook{} were Storybook
fixtures, throwaway tables such as \path{/vault} and \path{/sm/:organizationId} that exist only
to render one component in isolation; neither tool reads them now, and the oracle reports how
many files it set aside. \RoundTwoEarlierOracleBitwardenInside{} were inside the application.
\RoundTwoEarlierOracleBitwardenInsideVerbatim{} of those left the count for ghostfolio's
reason: they were recorded under a verbatim expression rather than an address ---
\path{/AuthRoute.Login} is what the earlier run stored, because the enum behind it was
unreadable --- so they leave that count and arrive in this one at the addresses they actually
mount. The rest are the mifos defect in miniature: the earlier run read their path segments
correctly but placed them at the project root, as \path{/password} and
\path{/session-timeout}, because it could not follow the lazy table that carries them; they are
declared in this run at the addresses the application mounts, under
\path{/settings/security/} and \path{/vault/}. Those are the same screens at a different
address, not screens gained or lost, which is a distinction a count of the denominator cannot
make and a diff of the set can. \RoundTwoEarlierOracleBitwardenAdded{} arrived, from two
causes: a route table wrapped in \texttt{Object.freeze} that both tools learned to read, and the
report and member subtrees, which were declared at the application root until
Section~\ref{sec:res-probe} showed the running router mounting them elsewhere. Against the
\RoundTwoEarlierOracleBitwardenNow{} screens that remain, the graph holds a node for every one.

An earlier draft reported bitwarden against 159 declared screens and decomposed
that figure in detail, and the decomposition was wrong in every part. The 159 came from a file
produced by a revision of the oracle from the middle of this round, dropped into the directory
the tables are generated from and read as though it were the earlier oracle's output. Nothing in
the generator could tell: the file is internally consistent, as every revision's output is, and
an absent file was read as an application whose denominator did not move, in silence. The
comparison is now made against the earlier oracle's own committed output, read from the
repository rather than from the directory the run writes into, so the revision each column
belongs to is a fact about the repository and not about what a working directory happened to
hold. The cost of not having done that is the paragraph this one replaces.

The two remaining moves are small and are printed for the same reason. alfresco went from
\RoundTwoEarlierOracleAlfrescoDeclared{} declared screens to
\RoundTwoEarlierOracleAlfrescoNow{}: \RoundTwoEarlierOracleAlfrescoInside{} left the count and
\RoundTwoEarlierOracleAlfrescoAdded{} arrived, because of the workspace-alias repair rather than
because anything was set aside. mifos went from
\RoundTwoEarlierOracleMifosDeclared{} to \RoundTwoEarlierOracleMifosNow{}, with
\RoundTwoEarlierOracleMifosInside{} paths leaving the count and
\RoundTwoEarlierOracleMifosAdded{} arriving: the routes of its reporting module, which the
earlier run placed at the project root as \path{/}, \path{/:filter} and \path{/run/:name}
because it would not follow a \texttt{baseUrl} import, now sit at \path{/reports} and beneath
it. Fictitious addresses left the set and real ones took their place, and because more left
than arrived the denominator moved down. The net change is too small to move any percentage in Chapter~\ref{ch:results},
which is exactly why it is printed: a denominator that moves without being mentioned is how the
previous figure survived three rounds of review, and the only defence that generalises is to
generate every part of the subtraction and let it be checked.

One caveat of Chapter~\ref{ch:evaluation} can now be priced. The oracle and the extractor agree
about an edge only where both read the same source, and \RoundTwoOutsideWitness{} edges across
\RoundTwoOutsideWitnessApps{} applications --- \RoundTwoOutsideWitnessWorstN{} of them in
\RoundTwoOutsideWitnessWorst{}, \RoundTwoOutsideWitnessWorstMust{} of those on the deterministic
tier --- have their only witness in a file outside the application directory. What the check
records is that the witness lies outside, not the route by which the file came to be read, so
the figure is reported here and not explained; the files are of the kind a monorepo puts beside
an application, a shared library and, in bitwarden, the repository's other front ends. Both tools read those files, so their agreeing there is not independent
evidence of anything; the witness check reports the figure rather than excluding the edges. The
reason is not that those edges are clicks in the application measured --- the files in
\RoundTwoOutsideWitnessWorst{} belong to a sibling front end built from the same repository, and
whether its components are reached in the application under measurement is not established here.
It is that a precision figure is a claim about the edges this tool emitted, wherever it read them,
and a reader who wants the figure over the measured directory alone can subtract the count. The
dead-link census of Section~\ref{sec:deadlinks} takes the other decision for the other reason:
calling a link dead asserts something about an application, and an accusation needs the link to be
reached in the one being measured, so there it is the boundary that decides and not the count.

The reason to distrust that last sentence is not in this application but in the mechanism, and
this round supplied the illustration. Three of bitwarden's declared screens were, for most of
this round, outside the denominator entirely: their paths are assembled inside a template
literal from an exported constant, neither tool could read them, and a denominator counted
without them gave this application a coverage of exactly 100\,\%. They were readable one
exported constant away. Both tools read them now and the figure is honest, but nothing about
the figure itself would have said otherwise: a screen neither tool can address is invisible to
a coverage ratio in exactly the way a screen that does not exist is. The only defence is to
report what was set aside, which is what the blind-spots column is for and why it is printed
beside every coverage figure in Chapter~\ref{ch:results} rather than in a footnote.

That so many reach every declared screen is a weaker result than it sounds, and the
blind-spots column is why: it counts the routes a tool knows are there and cannot address,
and reaching every screen one can see says the walk finds the route tables rather than that
the graph is complete. What the graph does with the navigation is the accounting below.


\section{Every remaining site names a reason}
\label{app:reasons}


\RoundTwoUnexplained{} sites remain unexplained, and Table~\ref{tab:round2reasons} gives the
reason for each. The reasons are produced by a script that exits non-zero if a single site
reaches the end of it without one, so a site the analysis cannot explain is a failing state
rather than a silent remainder. The largest group, \RoundTwoUnplacedClass{} of the
\RoundTwoUnexplained{}, is the attribution gap of Chapter~\ref{ch:analysis} showing through: a
class the extractor parsed and found a navigation site in, and could place under no screen,
because placing it needs the call graph with dependency injection that Section~\ref{sec:inputs}
specifies and this implementation does not build. The reason is the extractor's own note and not
this harness's inference, which matters because the alternative wording --- that nothing in the
application injects the class --- would be a claim about the application where the truth is a
claim about the tool. A second group concerns a target the extractor read and could not place: \RoundTwoTargetNoRoute{}
sites whose target matched no route, split by what can actually be claimed of each ---
\RoundTwoRelativeUnresolved{} a relative target with no base to resolve against,
\RoundTwoUnderUnreadable{} inside a subtree one of the tools could not read,
\RoundTwoNoRouteCaveated{} in an application whose route set is a lower bound, and
\RoundTwoDeadLinkUnanswered{} that nothing accounts for but the application. Each is decided
here by testing the target against the independent oracle's route set, pattern against
pattern, and not by repeating whatever note the extractor had left at the site: the note says
no node of the extractor's matched, which is a claim about the extractor.

\begin{table}[t]\small
\caption{Why each remaining site is not a transition. The script that produces this table
fails if any site has no reason.}
\label{tab:round2reasons}
\TabRoundTwoReasons
\end{table}


\section{The four dead-link corrections}
\label{app:deadlinkfixes}

Four corrections belong with this finding, and all four came from having it attacked. Three
of them are the same mistake made in three places, which is the more useful thing to report.

A target no route declares is not a target that goes nowhere. An application with a
\texttt{**} route catches every unmatched URL, and it does so in one of two ways: a
\texttt{redirectTo} forwards the URL to another screen, and a \texttt{component} renders one.
The extractor follows both, holds the resulting edge at may, and records a note saying the
target matched nothing and where the catch-all sent it. The third round followed only the
redirecting flavour --- which is not the flavour either application the surviving count comes
from uses, so the correction held for the one application it had been written against and
failed for the two that needed it.

Resolving the target is one claim and calling the link dead is another, and conflating them
is how the first version of this finding went wrong. bitwarden declares its authentication
paths through an enum the extractor had failed to read, so targets such as \texttt{/login}
matched nothing for a reason that was the extractor's, and the wildcard then forwarded every
one of them and raised that application's recovered figure by eight points.
That enum is read now, and so is the wrapper call that holds it, so those paths resolve
rather than being forwarded; what remains is the mechanism.
An earlier version of this
paragraph named three bitwarden paths assembled inside a template literal instead; those are
read now, which is what moved its denominator, and what is left there is a boundary rather
than a path. The note says which of the two it is, because a dropped subtree and an
unreadable path are not the same loss. The edge is drawn either way --- Angular really does
forward those URLs --- so what had to change was the claim attached to it, not the edge. The
extractor counts the declared paths it could not read and records one note per application
saying the route set is a lower bound, and that caveat is written into every wildcard note
from that application.

What the caveat then does to a candidate is the part this round changed. A single unreadable
declaration used to void every claim of this kind across the whole application, which is
severe and was deliberate: the alternative is a census that reports an extractor's blind spot
as an application's defect. It is also a harder question than the census asks. Whether an
application's whole route set was read is one question; whether \emph{this} target is served
is another, and the runtime check of Section~\ref{sec:res-probe} answers the second directly,
because the configuration it reads is the one the framework holds. Each caveated candidate is
therefore put to that configuration. \RoundTwoFrameworkMountsIt{} are withdrawn because the
framework serves a screen at the target, which makes them a measurement of what the extractor
failed to read rather than of what the application got wrong; \RoundTwoFrameworkPrefix{} are
withdrawn because the target is matched by a route that places no screen of its own, and
\RoundTwoFrameworkRedirects{} because it is answered by a redirect. \RoundTwoWildcardCaveated{}
sites keep the caveat. A matcher route leaves undecided only what it could claim --- the addresses
under its own prefix --- and reading it as leaving the whole application undecided withheld twelve
findings on a ground the harness's own navigation refutes, ghostfolio's matcher sitting at
\path{/admin} and reaching nothing above it. Where the configuration still declines, a navigation
is asked instead: ghostfolio's \path{/admin/users} is behind that matcher, and the recogniser
places a component there, so the three sites naming it are withdrawn rather than left undecided.
What survives is a candidate whose target the framework answers with nothing but its catch-all.

A route neither tool can read is indistinguishable, from a navigation site, from a route the
application never declares --- and this thesis credited the difference to the application
three times before the mechanism was general enough to stop it. A route array built by a
spread (\texttt{children: [\ldots adminRoutes]}) was not followed, so the routes it carries
were minted at the project root and links to them matched nothing. ghostfolio mounts its user
administration screen with Angular's \texttt{matcher} property rather than \texttt{path}:

\begin{quote}
\texttt{\{ component: GfAdminUsersComponent,}\\
\texttt{\ \ matcher: usersMatcher, \ldots \}}
\end{quote}

\noindent where \texttt{usersMatcher} consumes a URL whose first segment is \texttt{users}.
Neither tool read the \texttt{matcher} property at all, so \texttt{/admin/users} appeared in
no route set, and the previous round of this chapter published it as the one dead link it had
confirmed by hand --- the hand check having searched the application for \texttt{path:
'users'}, which is not how that route is spelled. alfresco declares
\texttt{nodes/:nodeId} and loads its children through a tsconfig \texttt{paths} alias, which
neither tool resolved, so \texttt{/nodes/:nodeId/rules} was absent from both route sets and
two live links to it were reported as dead. All three are now read. More to the point, both
tools now record the holes they cannot read --- a matcher route, an unresolvable lazy
boundary, a path expression that is not a literal --- and no claim that a path is undeclared
is counted inside one. The discipline is what generalises; the three fixes are not.

A prefix is not a target. The route-set test read only the first quoted segment of a commands
array, so \texttt{['/nodes', id, 'rules']} was tested as the path \texttt{/nodes}, which
alfresco declares no route at while \texttt{/nodes/:nodeId/rules} is exactly what it does
declare. The whole array is now read as a pattern, its computed segments as holes, and
matched against the declared routes pattern to pattern. That correction cuts both ways: it
cleared alfresco's two candidates and it found ten more sites in mifos that the first-segment
test had passed over.

What survives all four corrections is this. Of the \RoundTwoWildcardRescued{} sites a
catch-all route answers, \RoundTwoWildcardCaveated{} come from an application whose route set is
a lower bound and whose addresses no measurement can decide,
\RoundTwoFrameworkMountsIt{} have a target the framework's own configuration serves with a
screen --- a defect, but the extractor's --- \RoundTwoFrameworkPrefix{} a target it matches with
a route that places no screen, which could never have had a node, and
\RoundTwoFrameworkRedirects{} a target it answers with a redirect. \RoundTwoWildcardNoNode{} name
a route the application declares as a screen and the extractor minted no node for,
\RoundTwoWildcardPrefix{} name a route it declares as a \emph{path prefix},
\RoundTwoWildcardBlind{} address a subtree the oracle could not open, \RoundTwoWildcardPattern{}
are a pattern with a hole where the first segment should be and so name no single address,
\RoundTwoWildcardRelative{} are a relative target the analysis had no base to resolve,
\RoundTwoWildcardAppRoot{} name the application root, for \RoundTwoWildcardNoTarget{} the
record names no target at all, so nothing is established about where it points, and
\RoundTwoOutsideDirSites{} sit outside the directory the measurement is scoped to and are the
subject of the paragraph above; \RoundTwoDeadLinkWildcarded{} remain. Of the \RoundTwoTargetNoRoute{} unexplained sites whose
target named no route, \RoundTwoRelativeUnresolved{} are again a relative target with no base,
\RoundTwoUnderUnreadable{} address an unreadable subtree and \RoundTwoNoRouteCaveated{} come
from an application whose route set is a lower bound at all; \RoundTwoDeadLinkUnanswered{} remain.
Those two residuals are the dead-link candidates: \RoundTwoDeadLinkCandidates{} sites at
\RoundTwoDeadLinkLocations{} places in the source, \RoundTwoDeadLinkUrls{} distinct URLs, in
\RoundTwoDeadLinkApps{} applications --- \RoundTwoDeadLinkWhere{}. The count is of sites rather than of defects: one wrong
prefix written into a shared template is one mistake and several sites, and one site resolved
against several owning routes is one place to go and look and several targets. Every category
in the enumeration above except the last of each arm is a limit of one of the two tools rather
than a defect of an application, which is the proportion worth remembering about this kind of
finding: the residual is what is left after the tools stop accusing the applications of their
own blind spots.

One of those categories used to be unsettled, and is now measured. A route the application
declares as a path prefix answers no URL of its own unless it also declares an empty-path
child, and the previous round could not say which of the sites it excused on that ground
addressed a prefix the application really answers. The framework's own configuration says, per
address: of the \RoundTwoFrameworkPrefix{} sites now excused on that ground, every one has a
target the framework matches with a route that places no screen of its own. A navigation to
one of them, mifos's \path{/collections}, activates a last route carrying no component at all,
so the address answers with the application's chrome and an empty outlet --- which is what the
exclusion claimed all along, and is now a measurement rather than an inference.

It is a candidate list rather than a verdict, because the preceding four corrections are
evidence that the tools' blind spots are not yet all known, and a route both the extractor and
the oracle fail to read looks, from the outside, exactly like a route that does not exist.
What can be said is narrower than the previous round claimed. Each candidate was checked
against the application's own route files by hand, which is how the \texttt{appusers},
\texttt{manage-jobs} and \texttt{configurations} mismatches above are quoted; no candidate was
checked by navigating to it in a running application, which is the only check that would settle
one. Half of the objection has since been answered, and not by the route files. Every
candidate's target is put to the route configuration the framework itself holds, read out of
the running application by the check of Section~\ref{sec:res-probe}, so the premise a dead-link
claim rests on --- that the application mounts nothing at that address --- no longer rests on two
readings of the same files. The classification that premise uses was itself held against what
the framework does, on \ProbeAnswerChecked{} navigated addresses, of which it agrees with
\ProbeAnswerAgreed{}. That is what keeps the census from being empty: both applications it draws
from declare part of their route set unread, which by the rule above would suppress every
candidate they carry. What is still unchecked is the other premise, that the link is reached at
run time at all, and one step below it: a target is matched against the framework's patterns
rather than navigated to one address at a time, so a candidate's own URL has not been through
the recogniser. What the
method contributes is not the verdict but the list: a static transition graph puts every such
candidate on one page, with a file and a line each, instead of leaving them to be met at run
time.


\section{Asking the running framework}

\label{sec:res-probe}

Every result so far rests on two readings of the same route files. The witness check tests a
must edge's destination against the independent parser's route set, and Section~\ref{sec:protocol}
states the condition under which that test is vacuous: where both tools read a route table at
the wrong prefix, they agree on a wrong address and the check passes. Agreement between two
readers of one text is evidence only to the extent that the readers can differ. This section
reports a third reading that does not read the text at all.

\subsection{How the framework was asked}

Each application was built in its development configuration and loaded in a headless browser.
Once it had bootstrapped, the live \texttt{Router} was found through Angular's debug surface and
its route configuration read directly. Every lazily loaded subtree was expanded by calling the
application's own \texttt{loadChildren} and taking the routes out of whatever it resolved to: the
array itself for a standalone subtree, and the route providers of the module for one that loads
an \texttt{NgModule}. Nothing in the procedure parses a route file. The answer is the
configuration the framework holds, produced by the application's own code, so where it disagrees
with both static readers the disagreement cannot be an artefact of a reading they share.

No server was run. The method needs the application to boot, not to work: an address is counted
as mounted when the router's configuration places a component there on the primary outlet,
whether or not a guard would admit a visitor or a resolver would find data. All \ProbeApps{}
applications bootstrapped and \ProbeRoutes{} route entries were read, with \ProbeLazyFailures{}
lazily loaded subtrees left unexpanded and \ProbePageErrorApps{} applications raising an uncaught
error. \ProbeConsoleErrorApps{} logged errors to the console --- failed requests to the server
that is not there, for the most part --- which is expected and does not bear on the route
configuration, since the configuration is assembled before any of those requests is made.
\ProbeMatcherRoutes{} route in the corpus decides its address through a \texttt{UrlMatcher}; that
route and anything beneath it have no address any of the three readers can resolve to a string,
so they are outside both sides of the comparison rather than filed under the parent's address.
Table~\ref{tab:probe} reports the result.

Reading a configuration is itself a reading, and the one this comparison depends on is which of
its entries answer an address with a screen. That reading is checked against what the framework
does. The harness navigates a running application to an address and records the recogniser's
verdict --- the event Angular emits when a URL has matched, which fires before any guard runs, so
a guard sending every visitor to a sign-in page does not hide what was matched. Over
\ProbeAnswerApps{} applications and \ProbeAnswerAddresses{} addresses, chosen to include every
address this chapter names, the classification taken from the configuration agrees with the
navigation on \ProbeAnswerAgreed{} of the \ProbeAnswerChecked{} it decides, and declines to decide
\ProbeAnswerUndecided{}, which is the address the \texttt{UrlMatcher} route claims. The
comparison of Table~\ref{tab:probe} does not use that classification --- it decides for itself
which entries answer an address --- so the same navigations are held against its own list as well,
on the two implications that bind it: an address the framework serves with a screen it must hold,
and an address answered only by the catch-all it must not. That agrees on
\ProbeCmpAgreed{} of \ProbeCmpChecked{}. The addresses are those this chapter reasons about and
not a sample of the corpus: the \ProbeMissed{} screens the extractor does not mint are
\ProbeAnswerAddresses{} navigations away from being individually checked, and are not claimed to
have been. Two of
its answers are worth naming because the comparison turns on them. An address whose route carries
children places no screen of its own: a navigation to mifos's \path{/collections} activates a last
route with no component at all, so the address answers with the application's chrome and an empty
outlet. And an application with a catch-all answers every address, so the configuration mentioning
an address is no evidence that anything is mounted there --- in mifos a URL that matches nothing
reaches a not-found component, and in bitwarden and kubernetes/dashboard it is redirected to a
screen.

\subsection{What the framework mounts and the extractor does not}

The other direction is not symmetric with the first, and the thesis does not treat it as such. A
screen the extractor does not mint costs recall; the graph says less than it could. An address it
mints that nothing mounts would be a claim about an application that is false. The corpus holds
\ProbeMissed{} of the first and \ProbeInvented{} of the second.

Of the \ProbeMissed{}, \ProbeMissedGhostfolio{} are one construct in ghostfolio. Its product
comparison pages are not written down: a list of competitor records is mapped to routes and
spread into the table, and each path is a template string over a field of the record. The
extractor reads the table, mints the parent, and declares the spread a route source it could not
read.

\ProbeMissedMifos{} are in mifos, and they are one construct as well: its users module holds two
route tables, one for each identity provider the deployment may use, and chooses between them with
\texttt{!environment.OIDC.oidcServerEnabled}. That constant is not fixed by the sources. It is
computed from \texttt{window.env}, injected into the page at deployment time, so no reader of the
code can say which of the two tables the application mounts, and the extractor declares the
subtree unread. The two tables happen to declare the same five addresses with different
components, and reading either would have named those five addresses correctly --- which is what
the extractor used to do, by walking a table the module does not register under the prefix at all.
That it arrived at the right addresses was an accident of this application: the same reading gives
every prefix every table a file holds, and where two tables differ it mints the cross product.
The addresses are lost here and the loss is declared, which is the exchange this analysis is
built to make.

\ProbeMissedAlfresco{} are in alfresco and no two of them share a cause. Two are outside the
directory the measurement scopes to, though for different reasons: \path{/login} is declared by
the shell application that composes the measured library, and
\texttt{/view/\allowbreak{}authentication-\allowbreak{}confirmation} by an installed package
that shell provides. The second is outside it by a design choice the extractor states --- it does
not follow a dependency into a package directory, because a package declares no screen of the
application under measurement.

The third is inside it, and \path{/about} is not declared in TypeScript at all. The library ships \path{about.plugin.json},
whose single entry is

\begin{quote}\footnotesize\ttfamily
\{ "id": "app.about", "path": "about", "layout": "app.layout.main",\\
\hphantom{\{ }"component": "app.about.component" \}
\end{quote}

\noindent
and an extension registry binds the identifier \texttt{app.about.component} to a class at run
time. The route exists in data; no reader of the application's code can see it, and the
extractor's analysis is over the code. The probe sees it because the framework assembles the
table from the manifest before the probe reads the router. This construct and mifos's
runtime-selected table are both recorded in the approximations of
Section~\ref{sec:approximations}. alfresco's route set is a declared lower bound on other
grounds, so no claim of completeness is made over it either way.

The remaining \ProbeMissedBitwarden{} is bitwarden's vault, whose route is produced by a helper
that chooses between two components by feature flag.

One construct in ghostfolio therefore accounts for almost all of this column, and it is worth
saying what that means and what it does not. It does not mean the graph is missing those
screens by accident: the construct is one spread, the extractor says it could
not read it, and the \ProbeMissedGhostfolio{} pages behind it are generated entries that differ
only in a path segment and a data key. It does mean that the cost of a declared blind spot can be
far larger than the number of blind spots suggests, which is a thing a count of notes cannot
show and this comparison can.

\subsection{Nothing is missing that was not declared}

\ProbeMissedWithoutDeclaration{} applications hold a screen the framework mounts and the extractor
has no node for without the extractor having said that part of its route set was beyond it. Every
application in Table~\ref{tab:probe} that has a missing screen is an application that declared
one. The correspondence is per application and not per cause, which is the weaker of the two
readings and the only one measured here: alfresco declares that it could not read ten spreads,
and one of its \ProbeMissedAlfresco{} missing screens, \path{/about}, has nothing to do with a
spread --- it is declared in a JSON manifest and in no route file at all. What the measurement
supports is that no application is missing a screen while asserting that its route set is whole
--- not that every missing screen is behind the construct its application named. Until this
measurement the thesis could assert neither.

On this corpus the correspondence happens to hold in the other direction as well:
\ProbeDeclaredButComplete{} applications declare part of the route set unread and turn out to be
missing nothing, and \ProbeSilentAndComplete{} declare nothing and are missing nothing. That is
not a property of the analysis and is not offered as one. A declaration costs the route set
whatever the construct behind it happened to carry, and one that cost nothing is a tool being
careful rather than a tool being wrong; mifos's declaration cost nothing until this round, when
the extractor stopped reaching the addresses behind it through a table the module does not
register, and the same note then stood for five lost screens. The direction that matters is the
other one: a screen missing with no declaration would be a tool asserting a completeness it does
not have, and there are none of those. It also changes what the dead-link census of
Section~\ref{sec:deadlinks} can rest on. That census used to suppress every candidate from an
application whose route set is a declared lower bound, and a per-application licence was added to
lift the suppression where this comparison found the declared set complete after all. On the
figures above that licence lifts the suppression for \RoundTwoProbeLifted{} applications: both
applications the census draws from are missing a screen of their own --- mifos
\ProbeMissedMifos{} and bitwarden \ProbeMissedBitwarden{} --- so neither can be licensed. A licence of that shape asks a harder question than the census needs ---
whether a whole route set is complete --- and loses the census of any application that misses one
screen anywhere. Section~\ref{sec:deadlinks} therefore asks the question per candidate instead,
of the same framework configuration: a link whose target the framework serves nothing at is a dead
link whether or not the rest of the route set was read, and one whose target it does serve is not
a dead link however the extractor read it.

The correspondence did not hold when this check was first run, and it was not a near miss. Two
applications with unread route sources were asserting complete ones: alfresco, whose document
viewer is assembled by a helper call spread into the table, and bitwarden, whose vault route is
built the same way. Each spread was skipped without a word. Bitwarden's silence had already been
paid for, because with its caveat absent that application's sites entered the dead-link census ---
on a route set that was missing a screen. Both readers now record a spread they cannot follow as
a route source they did not read. Their rules for doing so are not the same: the extractor
declares the spread whose expression it cannot resolve to an array, the independent parser
declares every spread, since it follows none. The second is the more conservative of the two and
neither is silent, which is what the property requires.

\section{What a route table's evidence does not establish}
\label{sec:ownership}

One assumption underlies every approximation in the reader of route tables, and the others are
instances of it. A route table that no site in the application claims --- no
\texttt{loadChildren} names it, no spread or \texttt{children} carries it, no registration call
takes it --- is read as a table the application mounts at its root. The assumption is what lets
either tool find an application's routes however they are declared, and it is an
over-approximation: a table registered nowhere the compiled build reaches is mounted at no
address at all, and reading it at the root puts its entries into the graph as screens. It is
taken deliberately, because the alternative loses the routes of every application whose root
registration neither tool recognises, and the cases below are the places where it has been
narrowed rather than abandoned.

Three further approximations in that reader share one shape, and they are listed
together because each was found by attacking a repair of the previous one. All three are
declared in the tool's own notes; none was closed, and the reason in each case is that
closing it would refuse a construct the corpus actually uses.

The first concerns which file owns a lazily loaded route array. An Angular module loaded
through \texttt{loadChildren} usually keeps its routes in a separate routing module that it
imports, so the reader follows one import to find them. Electing that owner no longer trusts
the \emph{shape} of an array --- a menu or breadcrumb configuration carrying \texttt{path}
properties is never elected, and a file claimed beside an elected owner is reported as an
incomplete route set rather than dropped in silence. But a \texttt{Routes}-typed array is
evidence that its author declared it to some router, never that \emph{this} application mounts
it. A sibling application's typed table, reached through a workspace alias and used as data,
is therefore still elected and still minted under the lazy prefix, with a deterministic edge
into it. Requiring the elected owner to lie inside the application directory closes that
case and simultaneously refuses a lazily loaded route array that legitimately lives in a
shared library outside it, which is a pattern the corpus does use. No corpus application
distinguishes the two, so the gap is left open and declared rather than traded for a
construct that is known to occur.

The second is the same boundary one step earlier: the \texttt{loadChildren} specifier itself
is not bounded, so a loader whose import resolves into another project's alias root mounts
that project's table under this application's prefix. The obvious guard --- requiring the
loaded file to belong to the same project root as the route that loads it --- was applied
experimentally and refuses a lazily loaded module that one corpus application keeps in an
alias root, turning several of its declared screens into declared blind spots and lowering
its reported coverage. By path alone the analysis cannot tell a shared lazy library from
another application's route table, and that is the honest statement of the limit.

The third is a reading, not a defect. A routing module imported by two lazily loaded modules
and loaded by no route is read as a table of the project root rather than as children of both
importers. The alternative reading matches Angular more closely but lets one route adopt
another's table, which is how the first approximation above was reached; the conservative
reading is taken, and it costs the affected routes their prefix.

A fourth was found by asking what the test that keeps a lazy route from adopting the
application's root table actually read, and it is reported here as a repair rather than as a
limitation, because what the asking produced was a measurement no reading of the code would
have given. What the test read was the candidate file's own call to \texttt{forRoot} on
\texttt{RouterModule}, Angular's own discriminator for a root table, which is why it was sound
where it fired. It is also why it fired nowhere. An application that registers its root with
the standalone \texttt{provideRouter} API is not recognised at all, and two of the six have no
\texttt{forRoot} call anywhere; an application that declares its root array in one file and
registers it in another --- ghostfolio and kubernetes/dashboard both do --- is recognised in the
registering file, which holds no route, and not in the file a lazily loaded module could adopt.
Instrumenting the predicate made the limit exact rather than arguable: it was asked 698 times
across the six applications in the independent parser and 2798 times in the extractor, and
answered yes on none of them. A guard that never fires is worse than no guard, because a commit
message and a chapter credit it with a property nothing has shown it to have.

Both tools now mark root tables per \emph{array}, taken from the argument of either
registration API and resolved through the name it is nearly always written as, which is what
makes the question answerable in the shape applications actually write.

Several repairs of the same reader stand on constructed fixtures and on nothing else, and that
is worth saying plainly rather than leaving to be discovered. A file holding a menu beside an
imported routing module, a \texttt{forChild} given a table declared in another file, a
conditional resolved inside a registration call, a root table a spread elsewhere also carries,
and a wrapper that takes its own path as a parameter: each of these was reproduced, each
produced either an address the application does not mount or a declared loss where nothing was
lost, and each is now refused. None of the five occurs in this corpus. Every figure in
Chapter~\ref{ch:results} is therefore identical with and without them, so no result here is
evidence that they work --- their evidence is the fixtures, which fail when the rule each pins is
removed on its own. The same holds of the refusals in the independent parser --- the test that stops a lazy route
adopting a table another route owns is asked 174 times in the extractor and 3512 in the parser
and answers yes in neither --- so for those the corpus is evidence about the reading the two
tools share and not about the guards that bound it. The per-array root-table test is the
exception, and the distinction matters: it is asked a few hundred times and answers yes eight
times in each reader, and one of those answers is what mints alfresco's \path{/rules}. A rule
the corpus exercises and a rule only a fixture exercises are different kinds of claim, and this
chapter says which each is. The reverse also holds and is the reason for the paragraph: a corpus of six
applications is not a sample that makes an analysis sound, and five constructs it happens not to
contain were enough to break two of the properties this thesis claims. One further case came
with it: a table a feature library both hands to \texttt{provideRouter} and exports for lazy
loading is mounted at the root \emph{and} under the prefix that loads it, so a claim on its
file may not cover both mounts --- the root scan now mints a root-registered array whatever
claims its file, which is alfresco's \path{/rules}.

A fifth is the mirror of the first, and worth stating because it is where a conservative
refusal becomes an over-approximation. When the only site that loads a route table lies in a
test specification, the analysis declines the claim and scans the table from the project root,
so its routes are reported at the application's root address. The guard establishes that a site
the walk cannot reach may not justify suppressing the root scan; it does not establish that the
application mounts the table there. That makes it a further instance of the unclaimed
top-level-\texttt{Routes} over-approximation declared above, and it occurs nowhere in this
corpus.

\label{sec:named-outlet}
A sixth is the named outlet, and its history in this work is the clearest case of the difference
between declaring an approximation and being sound over it. A route carrying
\texttt{outlet: 'name'} is reached as \texttt{(name:path)}, so it contributes no screen on the
primary outlet. The extractor refused such a route from the start --- but only the route that
carries the property: where a named outlet opened a \texttt{children} subtree, the routes beneath
it were still assembled against the primary path and minted there. alfresco mounts its document
viewer that way, and the graph held \path{/view/:nodeId} and
\path{/search-libraries/view/:nodeId}, two addresses that application serves only through the
viewer outlet. The independent route parser read the subtree the same way, so coverage was
100\,\% over a set holding both, and no comparison between the two tools could show the gap. This
was written down as an open case, on the reasoning that an approximation stated is an
approximation a consumer can price.

That reasoning was wrong, and the runtime check of Section~\ref{sec:res-probe} is what showed it.
A consumer reading the graph finds a node at \path{/view/:nodeId}; a navigation to that address in
the built application is answered by the catch-all route. The declaration did not make the node
true, and the must-tier property admits no exception that a sentence can buy back. An outlet is
inherited, so nothing below a named-outlet route has a primary address either: the extractor now
carries outlet context down the walk and mints no screen for such a subtree. What remains is a
limitation on what the graph contains --- a screen reached only through a named outlet is not an
element of the node set, and the transitions into it are not in the edge set --- which is a
statement about recall and is the kind of approximation a declaration can honestly cover.

One attribution of a blind spot to the component it hides is a filename heuristic and is
labelled as one. Where a route prefix the analysis could not follow has literal segments, a
component whose path carries all of them as whole directory names is reported as hidden behind
that prefix. The heuristic answers the question a reader asks --- which screens are missing
because of this hole --- and it is wrong in two directions that the harness now states rather
than hides. A route prefix does not determine a directory name, so the components behind a hole
whose directory is named differently are not attributed to it and are reported only as having no
node; and a prefix made entirely of parameters has no literal segment, so it is attributed to no
component at all. An earlier version compared the prefix to the path as a bare substring, which
read \path{/cluster} into \texttt{resourcelist/clusterrole/component.ts} and would have replaced
the extractor's own true note about that file with a tool blind spot that does not cover it.
No count of the rows that satisfied the broader predicate is given, because the figure depends
on which population it is taken over and nothing in this thesis rests on it. What does matter
is why the attribution reaches the published reason table in neither form, and the obvious
account of that is wrong: it is not that some earlier test answers those rows first. The attribution is asked of
all \RoundTwoUnexplained{} unaccounted rows in the corpus and declines every one of them on its
own first two conditions --- \RoundTwoUnplacedClass{} because the owning file is not a component
and the rest because the graph already witnesses an edge from that file --- so the comparison
this paragraph is about is never reached on this corpus at all. What is recorded here is therefore a property of the rule and not
of any figure in Chapter~\ref{ch:results}. The narrower rule is still the one to state, because
the broader one was false wherever it was asked.
